\documentclass[11pt,a4paper,fontset=none]{ctexart}
\setCJKmainfont[Path=../../assets/fonts/,AutoFakeBold=2.5,AutoFakeSlant=.2]{NotoSansSC.ttf}
\setCJKsansfont[Path=../../assets/fonts/,AutoFakeBold=2.5,AutoFakeSlant=.2]{NotoSansSC.ttf}
\setCJKmonofont[Path=../../assets/fonts/,AutoFakeBold=2.5]{NotoSansSC.ttf}
\usepackage[margin=2.3cm]{geometry}
\usepackage{amsmath,amssymb}
\usepackage{booktabs,longtable,array}
\usepackage{xcolor}
\usepackage{listings}
\usepackage{graphicx}
\usepackage{fancyhdr}
\usepackage{hyperref}
\usepackage{enumitem}
\usepackage{tikz}
\usepackage{pgfplots}
\pgfplotsset{compat=1.18}
\hypersetup{colorlinks=true,linkcolor=blue!55!black,urlcolor=blue!65!black,citecolor=blue!55!black}
\urlstyle{same}
\setlength{\parindent}{2em}
\setlength{\parskip}{0.25em}
\setlength{\headheight}{14pt}
\setlength{\emergencystretch}{2em}
\pagestyle{fancy}
\fancyhf{}
\lhead{libsecp256k1：安全修复、性能优化与数学原理}
\rhead{\thepage}
\definecolor{labblue}{HTML}{2457A6}
\definecolor{labgreen}{HTML}{16825D}
\definecolor{laborange}{HTML}{C66A14}
\definecolor{codebg}{HTML}{F5F7FA}
\lstset{basicstyle=\ttfamily\small,backgroundcolor=\color{codebg},frame=single,breaklines=true,columns=fullflexible,showstringspaces=false}

\title{\bfseries libsecp256k1 的安全修复、性能优化与数学原理\\[0.5em]
\large 从“未绑定消息的伪造”到 constant-time 与 fixed-base multiplication}
\author{Bitcoin 实验一：任务二}
\date{2026 年 7 月 20 日}

\begin{document}
\maketitle

\begin{abstract}
Bitcoin 使用 secp256k1 椭圆曲线完成交易授权。其安全性不仅依赖 ECDSA/Schnorr 数学，也依赖消息摘要与交易语义的正确绑定、编译器不破坏 constant-time 代码，以及高效而可审计的有限域和点乘实现。本文以 \texttt{bitcoin-core/secp256k1} 稳定标签 \texttt{v0.7.1} 为基线，复核 \texttt{v0.3.1}、\texttt{v0.3.2} 的编译器相关 timing side-channel 修复，以及 \texttt{v0.4.1}、\texttt{v0.5.0/v0.5.1} 的性能演进。本文完整推导课件中的“验证器不检查被签消息”存在性伪造，并说明它攻击的是应用层消息绑定而非 ECDSA 的预定消息不可伪造性。七套 Release 构建的官方测试全部通过，其中新增 \texttt{v0.3.0/v0.3.1/v0.3.2} 三套安全版本构建；在 Apple ARM64/Clang 17 上，四套性能配置的 10 轮 \texttt{ecdsa\_sign} benchmark 中位数为 28.4、27.8、25.2 和 24.6 $\mu$s，展示 signed-digit multi-comb 算法与 2/22/86 KiB 预计算表的时间-空间权衡。

\textbf{关键词：} secp256k1；ECDSA；Schnorr；constant-time；timing side-channel；multi-comb；Bitcoin
\end{abstract}

\clearpage
\tableofcontents
\clearpage

\section{研究对象、证据与方法}

\texttt{libsecp256k1} 是面向 secp256k1 的高性能、高保证 C 库，提供 ECDSA、BIP340 Schnorr、ECDH、ElligatorSwift 与 MuSig2 等模块。本文采用三层证据：

\begin{enumerate}[leftmargin=2em]
  \item 官方 tag/commit/PR 与 CHANGELOG，固定“何时、为何、改了什么”；
  \item 源码结构和数学推导，解释安全或性能因果链；
  \item 同机同编译器 Release 构建、官方 ctest 和重复 benchmark，验证本机可复现部分。
\end{enumerate}

研究基线固定为表 \ref{tab:tag-commits} 所列 tag，而不是随时间移动的 master。脚本从本地官方 Git 仓库执行 \texttt{git rev-parse} 和 \texttt{git show} 生成 \texttt{tag\_commits.csv}；因此版本、commit 与日期不是人工抄写的常量。安全修复以修复前/后的 tag 构建，性能案例则以相同版本不同 CMake 参数做受控比较。

\begingroup
\setlength{\tabcolsep}{3pt}
\begin{longtable}{p{0.12\textwidth}p{0.21\textwidth}p{0.17\textwidth}p{0.40\textwidth}}
\caption{固定研究标签与 commit 对照}\label{tab:tag-commits}\\
\toprule
Tag & Commit & 日期 & 在本文中的用途 \\
\midrule
\endfirsthead
\toprule
Tag & Commit & 日期 & 在本文中的用途 \\
\midrule
\endhead
v0.3.0 & \texttt{bdf39000b9c6} & 2023-03-08 & Clang conditional-move 修复前基线，已 build/test \\
v0.3.1 & \texttt{346a053d4c44} & 2023-04-10 & 包含 PR \#1257，已 build/test \\
v0.3.2 & \texttt{acf5c55ae6a9} & 2023-05-13 & 包含 PR \#1303，已 build/test \\
v0.4.0 & \texttt{199d27cea322} & 2023-09-04 & v0.4 系列的源码差异边界 \\
v0.4.1 & \texttt{1ad5185cd42c} & 2023-12-21 & 移除 x86\_64 field assembly 后的 benchmark 基线 \\
v0.5.0 & \texttt{e3a885d42a78} & 2024-05-06 & multi-comb 与 2/22/86 KiB 三配置 \\
v0.5.1 & \texttt{642c885b6102} & 2024-08-01 & 默认预计算表调整与 API 文档证据 \\
v0.7.1 & \texttt{1a53f4961f33} & 2026-01-26 & 最新稳定研究基线与当前接口边界 \\
\bottomrule
\end{longtable}
\endgroup

\section{实验环境}

复现实验在同一台本地机器、同一编译器与 Release 配置下完成，环境信息由脚本写入 \texttt{results/environment.json}。关键配置如表 \ref{tab:environment}。

\begin{table}[htbp]
\centering
\caption{任务二本机实验环境}
\label{tab:environment}
\begin{tabular}{p{0.24\textwidth}p{0.66\textwidth}}
\toprule
项目 & 配置 \\
\midrule
硬件 & MacBook Pro（Mac14,9），Apple M2 Pro 10-core，16 GB memory \\
操作系统 & macOS 15.6.1，ARM64 \\
Python & CPython 3.13.3 \\
C 编译器 & Apple Clang 17.0.0，target \texttt{arm64-apple-darwin24.6.0} \\
构建工具 & CMake 4.4.0，Ninja 1.13.0 \\
构建配置 & \texttt{CMAKE\_BUILD\_TYPE=Release}；七套构建均启用 tests，四套性能配置启用 benchmark \\
测试 & CTest：\texttt{tests}、\texttt{noverify\_tests}、\texttt{exhaustive\_tests} \\
Benchmark & \texttt{SECP256K1\_BENCH\_ITERS=20000}；每组预热 2 次、正式运行 10 次 \\
大小统计 & 自动定位非符号链接的 \texttt{libsecp256k1} 动态库，记录 bytes 与 KiB \\
\bottomrule
\end{tabular}
\end{table}

平台边界必须明确：GCC 13.1 的 ECDH constant-time 回归和 x86\_64 handwritten assembly 的性能结论只能引用上游证据；本机 ARM64/Clang 数据不能冒充对 x86/GCC 结论的直接复现。为减少后台负载和温度变化带来的干扰，各配置采用相同构建类型、内部迭代次数、预热次数和正式运行次数；报告使用 median 与 IQR，而不是只报告单次最快结果。

\section{secp256k1 与 ECDSA 数学}

\subsection{有限域、曲线与群}

secp256k1 定义在素域 $\mathbb{F}_p$ 上：
\[
p=2^{256}-2^{32}-977,
\qquad E:\ y^2=x^3+7\pmod p.
\]
生成元 $G$ 的阶为素数 $n$：
\begin{lstlisting}[language={}]
n = FFFFFFFFFFFFFFFFFFFFFFFFFFFFFFFE
    BAAEDCE6AF48A03BBFD25E8CD0364141
\end{lstlisting}
私钥 $d\in\{1,\ldots,n-1\}$，公钥 $P=dG$。从 $P$ 求 $d$ 被认为等价于求解 elliptic-curve discrete logarithm problem（ECDLP）。域元素模 $p$，而签名中的 nonce、$r,s$ 等 scalar 模 $n$；混淆两种模数会造成实现错误。

对不同的仿射点 $P=(x_1,y_1)$、$Q=(x_2,y_2)$，群加法为
\begin{align*}
\lambda&=(y_2-y_1)(x_2-x_1)^{-1}\bmod p,\\
x_3&=\lambda^2-x_1-x_2\bmod p,\\
y_3&=\lambda(x_1-x_3)-y_1\bmod p.
\end{align*}
倍点时把斜率改为 $\lambda=3x_1^2(2y_1)^{-1}\bmod p$。这些公式说明 affine 坐标几乎每次加法都需要一次模逆；模逆明显比乘法昂贵。Jacobian 坐标用
\[
x=X/Z^2,\qquad y=Y/Z^3
\]
表示同一个仿射点，把点乘循环中的大量求逆替换为有限域乘法与平方，只在最终归一化时求逆。这里 field limb 是对模 $p$ 坐标的机器字分块，scalar limb 则是对模 $n$ 整数的分块；两套 reduction、边界与 constant-time 条件不能互换。

\subsection{ECDSA 签名与验证}

对消息 $m$，先定义 $e=H(m)$（按协议规则解释或截断）。签名者选择秘密 nonce $k\in\mathbb{Z}_n^*$：
\begin{align*}
R&=kG, & r&=x(R)\bmod n,\\
s&=k^{-1}(e+rd)\bmod n.
\end{align*}
签名为 $(r,s)$。验证者检查范围后计算
\begin{align*}
w&=s^{-1}\bmod n,\\
R'&=(ew)G+(rw)P,
\end{align*}
并接受当且仅当 $x(R')\bmod n=r$。正确性来自
\[
s^{-1}(eG+rP)
=k(e+rd)^{-1}(e+rd)G
=kG=R.
\]

nonce 不能复用：若同一 $k$ 对不同摘要 $e_1,e_2$ 产生 $(r,s_1),(r,s_2)$，则
\[
k=(e_1-e_2)(s_1-s_2)^{-1}\bmod n,
\quad d=(s_1k-e_1)r^{-1}\bmod n.
\]
库默认使用 RFC6979-style derandomized nonce，并对 secret-key 操作追求 constant-time 与 constant-memory-access。

\subsection{为什么验证可以 variable-time，而签名不可以}

签名计算 $kG$、私钥调整和 ECDH 点乘的 scalar 是秘密；分支方向、循环次数、表索引或 cache line 若依赖这些值，就可能形成 timing/cache/power side-channel。验签中的 $e,r,s$ 和公钥已经公开，使用 wNAF 或 variable-time modular inverse 通常不额外泄漏秘密。所谓“普通性能优化”只要求结果相同且平均运行更快；密码学 secret path 的优化还必须保持控制流、内存访问和运算轮数不依赖秘密。二者的威胁模型不同，不能仅凭同一套 functional tests 互相替代。

\section{课件中的“未绑定消息伪造”}

给定公钥 $P=dG$，攻击者任选 $u,v\in\mathbb{Z}_n^*$，计算
\[
R'=uG+vP=(x',y'),\qquad r'=x'\bmod n,
\]
并定义
\[
s'=r'v^{-1}\bmod n,
\qquad e'=r'uv^{-1}\bmod n.
\]
由于
\[
s'^{-1}=vr'^{-1},
\]
验证方得到
\begin{align*}
s'^{-1}(e'G+r'P)
&=(vr'^{-1})(r'uv^{-1}G+r'P)\\
&=uG+vP=R',
\end{align*}
因此 $(r',s')$ 对摘要 $e'$ 验证成功。这里并不知道私钥 $d$。

关键限制是：攻击者先得到随机式的 $e'$，并没有为预先指定的消息 $m$ 生成签名。若正确验证器重新计算并要求 $e'=H(m)$，攻击者还必须求一个满足该摘要的消息，这退化为对哈希函数的 preimage 问题。故课件展示的是：\textbf{如果应用只接受调用者传入的摘要，却不把摘要与实际消息重新绑定，就可产生存在性伪造。}它不是 Bitcoin ECDSA 已被破解。

这一边界也直接写在官方 API 文档中。\texttt{v0.5.1/include/secp256k1.h} 的 ECDSA verify 注释要求 verifier 自己对消息应用密码学哈希、不得直接接受外部 \texttt{msghash32}，否则无需知道私钥也容易构造“有效”签名。也就是说，课件推导不是脱离实现的理论提醒，而是库接口明确要求调用者承担的安全前置条件。

\section{Bitcoin 如何绑定被签内容}

Bitcoin 的签名不是对 explorer 展示文本签名，而是对严格序列化的 transaction digest 签名。不同脚本版本使用不同 sighash 规则：

\begin{itemize}[leftmargin=2em]
  \item Legacy ECDSA：从交易副本和 scriptCode 构造摘要；历史上计算复杂且容易出现二次哈希工作。
  \item SegWit v0/BIP143 ECDSA：显式承诺 prevouts、sequences、当前 outpoint、scriptCode、\textbf{输入金额}、outputs、locktime 与 sighash type。
  \item Taproot：BIP340 Schnorr 签名由 BIP341/342 的 tagged hash 绑定交易上下文、Taproot key/script path 等语义。
\end{itemize}

节点从交易和被花费输出自行重建 digest，不接受交易发送者随意声明“这个 hash 对应某消息”。因此上述 $e'$ 技巧不能把任意伪造签名附着到一笔指定的有效 Bitcoin 花费上。

Schnorr 采用形如 $sG=R+eP$ 的线性验证关系，便于 multisignature 与 batch verification；BIP340 通过 key prefixing、唯一 x-only 编码和 tagged hashing 防止相关构造中的歧义。它与 ECDSA 使用同一曲线，但签名方程和安全证明不同。

更具体地，BIP340 令 x-only 公钥为 $P$，nonce 点为 $R=kG$，并计算 tagged challenge
\[
e=H_{\mathrm{BIP0340/challenge}}(x(R)\parallel x(P)\parallel m)\bmod n,
\qquad s=k+ed\bmod n.
\]
验证检查 $sG=R+eP$，同时约束 $R$ 的偶数 $y$ 与坐标范围。challenge 同时包含 $R$、公钥 $P$ 与消息 $m$，所以线性方程本身不等于“消息未绑定”。BIP341 的 Taproot sighash 再把交易输入、输出、金额、scriptPubKey、annex 与 spend type 等上下文按规则纳入 $m$。

\begin{table}[htbp]
\centering
\setlength{\tabcolsep}{3pt}
\caption{secp256k1 在 Bitcoin 相关场景中的算法与秘密边界}
\label{tab:bitcoin-usage}
\begin{tabular}{p{0.20\textwidth}p{0.22\textwidth}p{0.24\textwidth}p{0.24\textwidth}}
\toprule
场景 & 算法/入口 & 被绑定的数据 & 性能/侧信道重点 \\
\midrule
Legacy 花费 & ECDSA + legacy sighash & 交易副本、scriptCode、sighash type & 签名 secret；验证 public \\
SegWit v0 & ECDSA + BIP143 & prevouts、金额、outputs 等 & 摘要可缓存，避免历史二次工作 \\
Taproot & BIP340 + BIP341/342 & tagged transaction digest 与 spend path & x-only key、Schnorr、batch-friendly \\
钱包密钥/签名 & key generation、ecmult\_gen & private key、nonce & 必须 constant-time、blinding \\
可选协议模块 & ECDH/MuSig 等 & 协议各自 transcript & scalar 常含秘密，不能套用公开验签假设 \\
\bottomrule
\end{tabular}
\end{table}

\section{安全修复案例}

\begingroup
\setlength{\tabcolsep}{3pt}
\begin{longtable}{>{\raggedright\arraybackslash}p{0.12\textwidth}>{\raggedright\arraybackslash}p{0.18\textwidth}>{\raggedright\arraybackslash}p{0.25\textwidth}>{\raggedright\arraybackslash}p{0.36\textwidth}}
\toprule
版本 & Merge commit / PR & 风险 & 修复要点 \\
\midrule
\endhead
v0.3.1 & \texttt{2d51a454}\newline PR \#1257 & Clang 15 能把源代码 cmov 重新编译成 branch & 在 field/scalar cmov 的 condition 上使用 volatile trick，阻止秘密相关控制流/内存访问 \\
v0.3.2 & \texttt{ab5a9171}\newline PR \#1303 & GCC 13.1 下 ECDH 路径可能出现 secret-dependent control flow & 在 scalar conditional negation 等更多“条件路径”谨慎使用 volatile，并扩展新编译器 CI \\
\bottomrule
\end{longtable}
\endgroup

\subsection{为什么“源代码无 if”仍不够}

常见 branch-free select 写成位掩码：
\[
z=(x\land \neg m)\lor(y\land m),
\]
其中 $m$ 由秘密条件扩展为全零或全一。编译器在抽象机语义上可以证明它等价于条件选择，并生成 branch 或 secret-indexed load；C 语言本身不承诺 timing behavior。故 constant-time 是“源码 + 编译器版本 + flags + 目标架构”的联合属性。

v0.3.1 的 PR \#1257 明确记录 Clang 15 对 field element cmov 的这种重写。将条件经 volatile 读取能限制某些优化，但它仍是工程防线而非语言级证明。因此项目随后增加对开发版 GCC/Clang、Valgrind/内存检查、无额外 VERIFY 的 production-like tests，并持续检查生成代码。

\subsection{v0.3.1：field/scalar cmov 的修复证据}

下面是固定 tag 中 \texttt{src/field\_5x52\_impl.h} 的关键差异；其余 limb 的 select 逻辑不变。
\begin{lstlisting}[language=C,caption={v0.3.0 到 v0.3.1 的条件读取变化}]
/* v0.3.0 */
mask0 = flag + ~((uint64_t)0);
mask1 = ~mask0;

/* v0.3.1 */
volatile int vflag = flag;
mask0 = vflag + ~((uint64_t)0);
mask1 = ~mask0;
\end{lstlisting}
当 flag 为 0 时，\texttt{mask0} 为全 1、保留原值；flag 为 1 时，\texttt{mask1} 为全 1、选择新值。问题不在这个代数掩码的功能正确性，而在优化器可能识别“二选一”语义并生成 secret-dependent branch 或 load。volatile 读取为该编译器优化建立屏障，使后续 bitwise select 保持预期形态。相同策略同步用于不同 field/scalar backend，而不是只修一个架构文件。

\subsection{v0.3.2：GCC 13.1 ECDH 路径的扩大修复}

v0.3.2 继续检查不叫 \texttt{cmov}、但同样由秘密 flag 控制的辅助函数。例如 \texttt{scalar\_cond\_negate} 从
\begin{lstlisting}[language=C,caption={v0.3.1 到 v0.3.2 的 conditional negate 变化}]
/* v0.3.1 */
uint64_t mask = !flag - 1;

/* v0.3.2 */
volatile int vflag = flag;
uint64_t mask = -vflag;
\end{lstlisting}
并在 \texttt{scalar\_cadd\_bit} 中同样通过 volatile 读取 flag。ECDH 入口先把外部 32-byte scalar 归约并以 conditional move 处理 overflow/zero，然后调用 constant-time \texttt{ecmult\_const}。因此 GCC 13.1 若在这些辅助条件上恢复秘密分支，最终可能污染 ECDH 的 constant-time 性质，即使输出点仍完全正确。修复的“why”是生成代码的时间行为而不是群运算公式错误；这也解释了为何 CTest 通过不能单独证明 side-channel 安全。

\subsection{影响边界}

timing side-channel 需要攻击者获得足够精细、重复的时间或微架构观测；它与“单次远程请求立即恢复私钥”不同。v0.3.2 的直接问题位于 optional ECDH module，不应写成 Bitcoin transaction verification 本身已泄漏私钥。但同一库被钱包、硬件设备、Lightning/P2P 组件使用，constant-time 回归仍属于必须快速升级的密码学工程风险。

\section{性能优化案例}

\subsection{域、标量与群运算}

库避免 floating point 和 runtime heap allocation。域元素模 $p$ 可用 5 个 52-bit limbs 或 10 个 26-bit limbs；scalar 模 $n$ 可用 4 个 64-bit 或 8 个 32-bit limbs。点通常使用 Jacobian $(X:Y:Z)$ 表示，使 affine $x=X/Z^2,y=Y/Z^3$，从而用乘法替代频繁且昂贵的 field inversion。

验证需要 $aP+bG$。库综合使用 wNAF、对固定生成元 $G$ 的大窗口预计算、Shamir's trick 和 secp256k1 endomorphism，将两个点乘合并并把部分 scalar 分裂为较短分量。模逆采用 safegcd 及其 variable-time 变体：secret-dependent 路径必须用 constant-time 版本；公开验证数据可使用 variable-time 版本换取速度。

wNAF 把 scalar 表示成稀疏的有符号奇数 digit，使非零项之间至少相隔若干位，从而减少点加法；Shamir's trick 在一次从高位到低位的扫描中联合计算 $aP+bG$，复用 doubling。secp256k1 还存在高效 endomorphism $\phi(P)=\lambda P$，可把 256-bit scalar 分解为两个约 128-bit 分量并行处理。fixed-base $kG$ 与 variable-base $kP$ 的区别在于 $G$ 永远固定，因而可以在编译期投入更大预计算表；任意 $P$ 则必须现场建立小表或使用不同算法。

\subsection{v0.4.1：删除 handwritten x86\_64 assembly}

PR \#1446（merge \texttt{10e6d29b}）删除 field operation 的手写 x86\_64 assembly。原因不是“assembly 天生慢”，而是现代编译器对对应 C 代码生成了更好的指令调度。官方 CHANGELOG 在 GCC 10.5 上报告 ECDSA verify 与 Schnorr verify 约 10\% 提升。该数字只适用于其 x86\_64/GCC 环境；本文 ARM64 benchmark 不把它当本机结果。

\subsection{v0.5.0：signed-digit multi-comb ecmult\_gen}

签名和公钥生成核心是 fixed-base multiplication $kG$。PR \#1058（merge \texttt{da515074}）引入 signed-digit multi-comb：将 scalar 重编码为 $\{-1,0,1\}$ 型 digit，把多个相隔固定 stride 的位组合为 table index，再以固定次数的 doubling、constant-time lookup 与 addition 完成点乘。实现同时处理：

\begin{itemize}[leftmargin=2em]
  \item 第一次 table lookup 不需要额外 point addition；
  \item 预计算避免不必要的 doublings；
  \item projective blinding 和未知离散对数的偏移点降低 side-channel 可控性；
  \item 提供 2、22、86 KiB 配置以适应嵌入式与桌面环境。
\end{itemize}

PR 上游记录 GCC 13.2 约 12.4\%、Clang 15 约 11.5\% 的提升。v0.5.1 的 PR \#1564（merge \texttt{ca06e58b}）将默认表从 22 KiB 调整到 86 KiB，与 Bitcoin Core 默认配置对齐。

\subsection{multi-comb 的代数推导与参数含义}

\texttt{v0.5.0/src/ecmult\_gen\_impl.h} 先定义 $B=\texttt{COMB\_BITS}$，以及
\[
\operatorname{comb}(s,P)=\sum_{i=0}^{B-1}(2s_i-1)2^iP
=\bigl(2s-(2^B-1)\bigr)P.
\]
为了计算 $R=g_nG$ 且避免模 2 除法，实现以随机 blinding scalar $b$ 设置
\begin{align*}
\textit{scalar\_offset}&=(2^B-1)/2-b\pmod n,\\
\textit{ge\_offset}&=bG,\\
d&=g_n+\textit{scalar\_offset}\pmod n,\\
R&=\operatorname{comb}(d,G/2)+\textit{ge\_offset}.
\end{align*}
代入第一式可见中间的 $bG$ 和 $(2^B-1)G/2$ 偏移最终严格抵消；随机化改变的是内部表示与中间点，不改变输出 $g_nG$。

若 block 编号为 $b$、每个 block 有 $T$ 个 teeth、spacing 为 $S$，源码定义
\[
\operatorname{mask}(b)=\sum_{t=0}^{T-1}2^{(bT+t)S},\qquad
\operatorname{table}(b,m)=\left(m-\frac{\operatorname{mask}(b)}{2}\right)G.
\]
最终计算可写成
\[
\sum_{i=0}^{S-1}2^i
\sum_{b=0}^{\mathrm{blocks}-1}
\operatorname{table}\!\left(b,(d\!\gg\! i)\land\operatorname{mask}(b)\right).
\]
外层从 $S-1$ 逆序到 0，每轮对所有 block 加一个表项，轮间倍点。表索引由秘密 scalar 决定，因此不能直接做 secret-indexed array load；实现遍历候选并用 cmov 选择。又因为翻转 mask 中所有 relevant bits 会得到相反点，只需存一半候选，表项数为
\[
\mathrm{blocks}\times 2^{\mathrm{teeth}-1}.
\]

表 \ref{tab:comb-parameters} 把三种 CMake 配置换算成源码注释给出的实际工作量。spacing 取 $\lceil256/(\mathrm{blocks}\times\mathrm{teeth})\rceil$，点加次数为 blocks$\times$spacing，倍点次数为 spacing$-1$。

\begin{table}[htbp]
\centering
\caption{v0.5.0 三种 ecmult\_gen 预计算配置的算法参数}
\label{tab:comb-parameters}
\begin{tabular}{lrrrrrr}
\toprule
配置 & blocks & teeth & spacing & 表项 & 点加 & 倍点 \\
\midrule
2 KiB  & 2  & 5 & 26 & 32   & 52 & 25 \\
22 KiB & 11 & 6 & 4  & 352  & 44 & 3 \\
86 KiB & 43 & 6 & 1  & 1,376 & 43 & 0 \\
\bottomrule
\end{tabular}
\end{table}

2 KiB 配置节省静态空间，但必须执行 25 次倍点；22 KiB 大幅减少 doubling，并把点加从 52 次降至 44 次；86 KiB 用 1,376 个预计算点把 spacing 降到 1，消除循环中的倍点，但 cmov 扫描和 cache footprint 更大，所以相对 22 KiB 的边际收益小于 2 到 22 KiB 的跃迁。这是可由源码参数解释的时间--空间权衡，而不只是 benchmark 表面相关性。算法来源可追溯到 Hamburg 的 compact elliptic-curve multiplication 工作\cite{hamburg}。

\section{本机复现实验}

\subsection{方法}

所有配置使用同一 Apple Clang 17、CMake Release、Ninja。安全案例增加 \texttt{v0.3.0}、\texttt{v0.3.1}、\texttt{v0.3.2} 三套 build，与原有四套性能 build 合计七套；每套均运行官方 \texttt{ctest}，三项测试（tests、noverify\_tests、exhaustive\_tests）全部通过。安全版本只做 build/test，避免混入性能对照；\texttt{ecdsa\_sign} 仍仅对四套性能配置各预热 2 次、正式运行 10 次，每次内部迭代 20,000，报告 10 次 Avg(us) 的 median 与 inclusive IQR。

每次正式运行产生一个独立观测 $t_i$。汇总脚本按
\[
\widetilde t=\operatorname{median}(t_1,\ldots,t_{10}),\qquad
\mathrm{IQR}=Q_{0.75}-Q_{0.25},\qquad
\mathrm{ops/s}=10^6/\widetilde t
\]
计算结果；相对基线的降时比例为 $(t_{\mathrm{base}}-t_{\mathrm{new}})/t_{\mathrm{base}}$。median 降低偶发调度尖峰的影响，IQR 展示中间 50\% 运行的离散度。预热结果不进入 CSV 的正式统计；所有原始行保留，报告数字可以从 \texttt{benchmark\_runs.csv} 重新生成。

\begin{table}[htbp]
\centering
\caption{Apple ARM64/Clang 17 的 Release 构建与测试矩阵}
\begin{tabular}{lll}
\toprule
配置 & 用途 & CTest \\
\midrule
v0.3.0 default & Clang constant-time 修复前基线 & 3/3 passed \\
v0.3.1 default & Clang 修复后、GCC ECDH 修复前 & 3/3 passed \\
v0.3.2 default & GCC ECDH constant-time 修复后 & 3/3 passed \\
v0.4.1 default & 性能对照基线 & 3/3 passed \\
v0.5.0 2 KiB & 小预计算表 & 3/3 passed \\
v0.5.0 22 KiB & 中预计算表 & 3/3 passed \\
v0.5.0 86 KiB & 大预计算表 & 3/3 passed \\
\bottomrule
\end{tabular}
\end{table}

\begin{table}[htbp]
\centering
\caption{Apple ARM64/Clang 17 的 ECDSA signing 结果}
\begin{tabular}{lrrrr}
\toprule
配置 & Median ($\mu$s) & IQR ($\mu$s) & ops/s & dylib bytes \\
\midrule
v0.4.1 default & 28.4 & 0.25 & 35,211 & 1,336,120 \\
v0.5.0 2 KiB & 27.8 & 0.15 & 35,971 & 1,270,344 \\
v0.5.0 22 KiB & 25.2 & 0.20 & 39,683 & 1,286,856 \\
v0.5.0 86 KiB & 24.6 & 0.10 & 40,650 & 1,352,904 \\
\bottomrule
\end{tabular}
\end{table}

\begin{figure}[htbp]
\centering
\begin{tikzpicture}
\begin{axis}[
  ybar,
  width=0.88\textwidth,
  height=6.2cm,
  ymin=0,
  ymax=32,
  ylabel={Median $\mu$s/signature（越低越好）},
  symbolic x coords={v0.4.1,2KiB,22KiB,86KiB},
  xtick=data,
  nodes near coords,
  bar width=18pt,
  enlarge x limits=0.18,
  ymajorgrids=true,
  grid style={dashed,gray!30}
]
\addplot[fill=labblue!70] coordinates {(v0.4.1,28.4) (2KiB,27.8) (22KiB,25.2) (86KiB,24.6)};
\end{axis}
\end{tikzpicture}
\caption{不同 fixed-base multiplication 配置的本机签名时间}
\end{figure}

相对 v0.4.1，2/22/86 KiB 配置的时间分别下降 2.11\%、11.27\%、13.38\%。同一 v0.5.0 内，22 KiB 相比 2 KiB 增加约 16.1 KiB dylib 大小并显著降时；86 KiB 再增加约 64.5 KiB，收益缩小，体现 cache footprint 与 lookup/addition 次数的权衡。

不同 tag 的 dylib 总大小还受版本代码变化影响，不能把 v0.4.1 与 v0.5.0 的总大小差异完全归因于表；只有 v0.5.0 三配置之间是受控比较。类似地，本机 11.27\% 与上游 Clang 数据接近，但架构、编译器版本和 CPU 不同，属于相互一致的独立观察，不是逐位复现。

library size 也由同一个 \texttt{reproduce.py} 自动生成：脚本在每个 build tree 中排除符号链接，定位实际 \texttt{libsecp256k1} 动态库并写入 bytes/KiB、tag 和 table 配置。这样删除 \texttt{work/builds} 后执行一次 \texttt{--case all}，能重新产生 build/test 日志、逐次 benchmark、统计汇总和 \texttt{library\_sizes.csv}，避免报告中出现不可追溯的手填大小。

\section{正确性、性能与 constant-time 的关系}

三项目标不能互相替代：

\begin{description}[style=nextline]
  \item[功能正确性] 对合法输入给出正确签名/验证结果；官方 vectors、exhaustive tests 和 Wycheproof 覆盖边界。
  \item[密码学安全性] nonce、消息绑定、输入范围和群运算满足方案假设，无法用代数关系伪造预定消息。
  \item[侧信道安全性] 执行路径、内存访问和功耗不泄漏 secret；即使数学输出正确，编译器引入 branch 仍可能失败。
\end{description}

性能优化必须保持前两项。\texttt{libsecp256k1} 的策略是固定次数运算、constant-time table lookup、branch-free conditional move、blinding、无堆分配、严格 API 与广泛测试。公开数据验证可以安全采用 variable-time 算法，而 signing/private-key generation 不能因为 benchmark 更快就接受 secret-dependent 路径。

\subsection{实验能证明什么，不能证明什么}

七套构建与 CTest 证明这些 tag 在本机 Release 配置下可以编译，且官方功能/穷举测试通过；四组 benchmark 量化 Apple ARM64/Clang 17 上特定签名微基准的相对性能。它们不能证明任意编译器生成代码均 constant-time，也不能把微基准收益直接等同于钱包端到端吞吐量。本文没有在 x86\_64 上重跑 GCC 13.1 或 GCC 10.5，故 v0.3.2 与 v0.4.1 的平台特有结论只依据官方 PR/CHANGELOG；本机 build/test 是版本可构建性的补充证据，而非该侧信道或 assembly 性能结果的直接复现。

此外，macOS 动态调频、温度、后台调度与 cache 状态仍可能影响微秒级测量。相同机器、相同工具链、预热、重复运行和 median/IQR 只能降低而不能消除这些因素。三种 v0.5.0 表大小之间的受控比较最强；跨 tag 的总库大小与运行时间还叠加其他源代码变化，解释时应保持这一因果边界。

\section{结论}

课件的伪造推导揭示了一个通用原则：签名验证必须重新建立“消息语义 $\rightarrow$ 规范序列化 $\rightarrow$ digest $\rightarrow$ signature”的完整链条。Bitcoin 的 sighash 正是该绑定，而不是让交易发送者自由声明 digest。

\texttt{v0.3.1/v0.3.2} 表明 constant-time 不是看源码有没有 \texttt{if}，而是需要持续对新编译器与目标架构审计。\texttt{v0.4.1/v0.5.x} 表明高级数学表示、现代编译器和预计算表共同决定性能。本机数据验证了 signed-digit multi-comb 在 ARM64 上的收益，并量化了 2/22/86 KiB 的时间-空间折衷；同时严格保留了 x86/GCC 特有结论的证据边界。

\clearpage
\appendix
\section{复现命令与数据位置}

\begin{lstlisting}[language=bash]
python3 -m venv .venv
.venv/bin/python -m pip install cmake ninja
PATH="$PWD/.venv/bin:$PATH" \
  .venv/bin/python scripts/reproduce.py \
  --case all --runs 10 --warmups 2
\end{lstlisting}

逐次数据位于 \texttt{results/benchmark\_runs.csv}，汇总位于 \texttt{results/benchmark\_summary.csv}；动态库与预计算表配置的大小位于 \texttt{results/library\_sizes.csv}。环境、tag commit、case commit、build/test 原始日志和版本区间 diffstat 均保存在 \texttt{results/}。上游源码与 build tree 位于被忽略的 \texttt{work/}，删除后可由脚本重建。若已有 build tree，可运行 \texttt{--case sizes} 单独重建库大小 CSV/JSON；\texttt{--case all} 会在构建后自动生成它们。

\begin{thebibliography}{20}
\bibitem{secpreadme} Bitcoin Core, \emph{libsecp256k1 README}. \url{https://github.com/bitcoin-core/secp256k1}
\bibitem{changelog} Bitcoin Core, \emph{libsecp256k1 CHANGELOG}. \url{https://github.com/bitcoin-core/secp256k1/blob/master/CHANGELOG.md}
\bibitem{pr1257} PR \#1257, \emph{ct: Use volatile trick in all fe/scalar cmov implementations}. \url{https://github.com/bitcoin-core/secp256k1/pull/1257}
\bibitem{pr1303} PR \#1303, \emph{ct: Use more volatile}. \url{https://github.com/bitcoin-core/secp256k1/pull/1303}
\bibitem{pr1446} PR \#1446, \emph{field: Remove x86\_64 asm}. \url{https://github.com/bitcoin-core/secp256k1/pull/1446}
\bibitem{pr1058} PR \#1058, \emph{Signed-digit multi-comb ecmult\_gen algorithm}. \url{https://github.com/bitcoin-core/secp256k1/pull/1058}
\bibitem{pr1564} PR \#1564, \emph{Adjust default precomputed signing table}. \url{https://github.com/bitcoin-core/secp256k1/pull/1564}
\bibitem{hamburg} M. Hamburg, \emph{Fast and compact elliptic-curve cryptography}, IACR Cryptology ePrint Archive, Report 2012/309. \url{https://eprint.iacr.org/2012/309}
\bibitem{bip143} BIP 143, \emph{Transaction Signature Verification for Version 0 Witness Program}. \url{https://github.com/bitcoin/bips/blob/master/bip-0143.mediawiki}
\bibitem{bip340} BIP 340, \emph{Schnorr Signatures for secp256k1}. \url{https://github.com/bitcoin/bips/blob/master/bip-0340.mediawiki}
\bibitem{bip341} BIP 341, \emph{Taproot}. \url{https://github.com/bitcoin/bips/blob/master/bip-0341.mediawiki}
\bibitem{safegcd} D. J. Bernstein and B.-Y. Yang, \emph{Fast constant-time gcd computation and modular inversion}. \url{https://gcd.cr.yp.to/}
\end{thebibliography}

\end{document}
